\section{A sharp joint range and phase countermodels}
\label{sec:sharp-range}

The factorization gives the upper edge
\(\Flat\le\Seff\).  There is also a universal lower edge when the
magnitude profile is too concentrated to cancel completely.

\begin{theorem}[Sharp effective-support strip]
\label{thm:sharp-strip}
For every nonzero finitely supported complex sequence,
\begin{equation}
 \boxed{
 (2-\Seff)_+\le\Flat\le\Seff.}
\label{eq:sharp-strip}
\end{equation}
The upper equality is characterized by
\cref{thm:equality-dictionary}.  If \(\Seff<2\), equality on the
lower edge occurs exactly in the following cases:
\begin{enumerate}[label=(\alph*),leftmargin=2.2em]
\item \(S=1\); or
\item \(S=2\) and the two coefficients are antiparallel.
\end{enumerate}
If \(\Seff\ge2\), equality on the displayed zero lower edge occurs
exactly when \(Z=0\).
\end{theorem}

\begin{proof}
Writing the support as a finite ordered set,
\begin{align*}
 |Z|^2
 &=
 D+2\sum_{m<n}\operatorname{Re}(a_m\overline{a_n})\\
 &\ge
 D-2\sum_{m<n}|a_m||a_n|\\
 &=
 D-(M^2-D)
 =
 2D-M^2.
\end{align*}
Division by \(D\) gives
\(\Flat\ge2-\Seff\); nonnegativity gives the positive part.
The upper edge is already in \cref{thm:exact-factorization}.

When \(\Seff<2\), the lower right side is positive.  Equality requires
\[
 \operatorname{Re}(a_m\overline{a_n})
 =-|a_m||a_n|
\]
for every pair of nonzero coefficients.  Thus every pair must be
antiparallel.  This is possible for one coefficient, vacuously, or
for two coefficients, and is impossible for three nonzero complex
numbers.  Conversely, one coefficient or two antiparallel
coefficients gives equality.  When \(\Seff\ge2\), the displayed
lower edge is zero, so equality there is equivalent to
\(\Flat=0\), hence to \(Z=0\).
\end{proof}

\begin{proposition}[Sharpness of the lower edge up to its zero endpoint]
\label{prop:positive-edge-sharp}
Every value \(s_{\rm e}\in[1,2]\) occurs on the lower edge
of \eqref{eq:sharp-strip}; the endpoint \(s_{\rm e}=2\) is the zero
endpoint.  More precisely, \(s_{\rm e}=1\)
is attained by one nonzero coefficient, while for every
\(1<s_{\rm e}\le2\) there are two nonzero antiparallel coefficients
for which
\[
 \Seff=s_{\rm e},
 \qquad
 \Flat=2-s_{\rm e}.
\]
\end{proposition}

\begin{proof}
For \(s_{\rm e}=1\), take the one-point sequence \(a=(1)\).
For the remaining values, take \(a=(1,-t)\) with \(0<t\le1\).
Then
\[
 \Seff(t)=\frac{(1+t)^2}{1+t^2},
 \qquad
 \Flat(t)=\frac{(1-t)^2}{1+t^2}
 =2-\Seff(t).
\]
The continuous function \(\Seff(t)\) increases from \(1\) to \(2\).
\end{proof}

\begin{proposition}[Pointwise sharpness of both edges]
\label{prop:both-edges-sharp}
For every \(s_{\rm e}\ge1\), some nonzero finite sequence satisfies
\[
 \Seff=s_{\rm e},
 \qquad
 \Flat=\Seff.
\]
For every \(s_{\rm e}\ge2\), some nonzero finite sequence satisfies
\[
 \Seff=s_{\rm e},
 \qquad
 \Flat=0.
\]
Thus both pieces of the boundary in \eqref{eq:sharp-strip} are
pointwise attainable.
\end{proposition}

\begin{proof}
Choose an integer \(k\ge s_{\rm e}\).  For one coefficient of
magnitude \(1\) and \(k-1\) coefficients of common magnitude
\(t\in[0,1]\), all with the same phase, the effective support is
\[
 E_k(t)
 =
 \frac{(1+(k-1)t)^2}{1+(k-1)t^2}.
\]
It increases continuously from \(1\) to \(k\), and the common-phase
condition gives \(\Flat=\Seff\).

For the zero edge, choose \(k\) with \(2k\ge s_{\rm e}\) and use
opposite pairs with magnitudes
\[
 (1,1,t,t,\ldots,t,t).
\]
Their signed coefficients may be written
\[
 (1,-1,t,-t,\ldots,t,-t).
\]
They sum to zero, while their effective support is
\[
 2E_k(t),
\]
which increases continuously from \(2\) to \(2k\).  Hence every
\(s_{\rm e}\ge2\) occurs with \(\Flat=0\).
\end{proof}

\begin{example}[Same magnitude level, opposite phase outcomes]
\label{ex:same-magnitude-opposite-phase}
Let \(S\ge2\) be even and \(L>0\).  On a support of size \(S\), let
\[
 a_n=L
\]
for every \(n\), and let \(b_n\) consist of \(S/2\) values \(L\)
and \(S/2\) values \(-L\).  Both sequences have
\[
 M=SL,\qquad D=SL^2,\qquad \Seff=S.
\]
But
\[
 \canc(a)=1,\quad \Flat(a)=S,
 \qquad
 \canc(b)=0,\quad \Flat(b)=0.
\]
Thus even exact natural effective support does not imply any phase
cancellation.
\end{example}

\begin{remark}[Complex phases cause no hidden hypothesis]
All proofs are over \(\C\).  The terms ``common phase'' and
``antiparallel'' are equality conditions in the complex triangle
inequality, not assumptions that the physical coefficient is real.
\end{remark}
